\section{Introduction et contexte}
\label{sec:intro}

Le marketing digital occupe aujourd'hui une place centrale dans la stratégie
commerciale des entreprises, qu'il s'agisse de pure-players e-commerce, de
marques traditionnelles en transition numérique ou d'agences spécialisées. La
multiplication des canaux d'acquisition (Meta, TikTok, Google Ads, e-mail,
SEO\ldots), des formats publicitaires (vidéo, image, carrousel, UGC) et des
indicateurs de performance (CPM, CPA, ROAS, LTV) a transformé le métier de
\emph{marketeur} en celui d'un véritable \emph{ingénieur de données}.

Pourtant, la grande majorité des équipes marketing continue d'opérer à l'aide
d'un patchwork d'outils mal intégrés~: Trello pour les idées de produits,
Google Sheets pour le suivi budgétaire, Notion pour les briefs créatifs,
Google Drive pour les assets vidéo, et enfin les interfaces natives des régies
publicitaires pour le pilotage des campagnes. Cette fragmentation engendre
trois problèmes structurels~: \textbf{la perte de connaissance}
(les angles psychologiques qui ont fonctionné disparaissent dans les fils Slack),
\textbf{l'inefficacité opérationnelle} (le calcul manuel du ROI consomme un
temps précieux), et \textbf{l'absence d'aide à la décision en temps réel}
(les équipes lancent de nouveaux produits sans appui analytique).

Le présent projet répond à cette problématique en proposant un \textbf{CRM
spécialisé pour le marketing de performance}~: une plateforme web unifiée qui
centralise le pipeline produit, la bibliothèque d'assets (angles, créatifs,
hooks), le suivi budgétaire et l'analyse du ROI, le tout enrichi par un
assistant conversationnel propulsé par l'API DeepSeek et par un modèle
prédictif d'évaluation du potentiel d'un nouveau produit.

\subsection{Objectifs du projet}

\begin{itemize}
    \item Concevoir et développer une application web Django couvrant
    l'intégralité du cycle de vie d'une campagne marketing.
    \item Garantir une expérience utilisateur multilingue native en
    \textbf{anglais, français et arabe} (avec gestion correcte du sens
    d'écriture droite-à-gauche).
    \item Intégrer un \textbf{plugin chatbot} basé sur l'API DeepSeek qui
    apporte une valeur métier concrète~: répondre aux questions du marketeur
    en s'appuyant sur les données réelles du CRM.
    \item Ajouter une fonctionnalité bonus de \textbf{prédiction par
    apprentissage automatique} permettant de catégoriser le potentiel
    (Faible / Moyen / Élevé) d'un nouveau produit.
    \item Produire l'ensemble des livrables académiques requis~: cahier des
    charges, diagrammes UML, présentation de soutenance.
\end{itemize}

\subsection{Public cible}

Trois profils d'utilisateurs ont été identifiés~:
\begin{description}
    \item[Le marketing manager] qui pilote les campagnes au quotidien et
    consulte le pipeline pour décider des prochains lancements.
    \item[L'analyste data] qui consulte les tableaux de bord ROI et utilise
    le chatbot pour interroger les données sans écrire de SQL.
    \item[L'administrateur] qui gère les utilisateurs, les permissions et
    les paramètres globaux.
\end{description}
